\documentclass{amsart}
% For dealing with references we use the comment environment
\usepackage{verbatim}
\newenvironment{reference}{\comment}{\endcomment}
%\newenvironment{reference}{}{}
\newenvironment{slogan}{\comment}{\endcomment}
\newenvironment{history}{\comment}{\endcomment}

% For commutative diagrams we use Xy-pic
\usepackage[all]{xy}

% We use 2cell for 2-commutative diagrams.
\xyoption{2cell}
\UseAllTwocells

% We use multicol for the list of chapters between chapters
\usepackage{multicol}

% This is generally recommended for better output
\usepackage{lmodern}
\usepackage[T1]{fontenc}

% For cross-file-references

% Package for hypertext links:
\usepackage{hyperref}

% For any local file, say "hello.tex" you want to link to please

% Theorem environments.
%
\theoremstyle{plain}
\newtheorem{theorem}[subsection]{Theorem}
\newtheorem{proposition}[subsection]{Proposition}
\newtheorem{lemma}[subsection]{Lemma}

\theoremstyle{definition}
\newtheorem{definition}[subsection]{Definition}
\newtheorem{example}[subsection]{Example}
\newtheorem{exercise}[subsection]{Exercise}
\newtheorem{situation}[subsection]{Situation}

\theoremstyle{remark}
\newtheorem{remark}[subsection]{Remark}
\newtheorem{remarks}[subsection]{Remarks}

\numberwithin{equation}{subsection}

% Macros
%
\def\lim{\mathop{\mathrm{lim}}\nolimits}
\def\colim{\mathop{\mathrm{colim}}\nolimits}
\def\Spec{\mathop{\mathrm{Spec}}}
\def\Hom{\mathop{\mathrm{Hom}}\nolimits}
\def\Ext{\mathop{\mathrm{Ext}}\nolimits}
\def\SheafHom{\mathop{\mathcal{H}\!\mathit{om}}\nolimits}
\def\SheafExt{\mathop{\mathcal{E}\!\mathit{xt}}\nolimits}
\def\Sch{\mathit{Sch}}
\def\Mor{\mathop{\mathrm{Mor}}\nolimits}
\def\Ob{\mathop{\mathrm{Ob}}\nolimits}
\def\Sh{\mathop{\mathit{Sh}}\nolimits}
\def\NL{\mathop{N\!L}\nolimits}
\def\CH{\mathop{\mathrm{CH}}\nolimits}
\def\proetale{{pro\text{-}\acute{e}tale}}
\def\etale{{\acute{e}tale}}
\def\QCoh{\mathit{QCoh}}
\def\Ker{\mathop{\mathrm{Ker}}}
\def\Im{\mathop{\mathrm{Im}}}
\def\Coker{\mathop{\mathrm{Coker}}}
\def\Coim{\mathop{\mathrm{Coim}}}

% Boxtimes
%
\DeclareMathSymbol{\boxtimes}{\mathbin}{AMSa}{"02}

%
% Macros for moduli stacks/spaces
%
\def\QCohstack{\mathcal{QC}\!\mathit{oh}}
\def\Cohstack{\mathcal{C}\!\mathit{oh}}
\def\Spacesstack{\mathcal{S}\!\mathit{paces}}
\def\Quotfunctor{\mathrm{Quot}}
\def\Hilbfunctor{\mathrm{Hilb}}
\def\Curvesstack{\mathcal{C}\!\mathit{urves}}
\def\Polarizedstack{\mathcal{P}\!\mathit{olarized}}
\def\Complexesstack{\mathcal{C}\!\mathit{omplexes}}
% \Pic is the operator that assigns to X its picard group, usage \Pic(X)
% \Picardstack_{X/B} denotes the Picard stack of X over B
% \Picardfunctor_{X/B} denotes the Picard functor of X over B
\def\Pic{\mathop{\mathrm{Pic}}\nolimits}
\def\Picardstack{\mathcal{P}\!\mathit{ic}}
\def\Picardfunctor{\mathrm{Pic}}
\def\Deformationcategory{\mathcal{D}\!\mathit{ef}}

\setlength{\emergencystretch}{3em}
\begin{document}
\title[Stacks Project, Tag 01ZM]{1667 --- Finitely presented schemes, morphisms and equality descend through inverse systems of quasi-compact quasi-separated bases with affine transition maps}
\maketitle
\noindent Copyright (C) 2005--2025 Johan de Jong. Stacks Project authors. Source: \href{https://stacks.math.columbia.edu/tag/01ZM}{Tag 01ZM}.

\noindent Editorial difficulty note (not author text). Difficulty H2: This categorical finite-presentation descent theorem must descend entire scheme objects, all morphisms and equality relations across arbitrary affine-transition inverse systems. Its proof constructs finite affine models, synchronizes overlap opens and isomorphisms at successive indices, and enforces the cocycle before gluing, without a Noetherian or separatedness assumption. The global categorical coherence construction is beyond routine ring coefficient descent..

\noindent Editorial provenance note (not author text). History: excerpt from revision \texttt{a0444\allowbreak 6e57e\allowbreak c1fbc\allowbreak 252a8\allowbreak 71afc\allowbreak ec775\allowbreak 2fb28\allowbreak 07b14}, limits.tex, lines 2369--2545. Reference presentation was adapted; original-author context and core support blocks were retained or added. The mathematical wording is unchanged. Licensed under GNU FDL 1.2 or later; no invariant sections or cover texts. See ../licenses/COPYING. Context blocks reproduce author definitions and supporting results; they are not additional counted problems.

\begin{situation}
\label{situation-descent}
Let $S = \lim_{i \in I} S_i$ be the limit of a directed system of schemes
with affine transition morphisms $f_{i'i} : S_{i'} \to S_i$
(Lemma \ref{lemma-directed-inverse-system-has-limit}).
We assume that $S_i$ is quasi-compact and quasi-separated for all $i \in I$.
We denote $f_i : S \to S_i$ the projection.
We also choose an element $0 \in I$.
\end{situation}

\begin{definition}
\label{properties-definition-ample}
\begin{reference}
\href{https://stacks.math.columbia.edu/bibliography}{Bibliography: EGA (II Definition 4.5.3)}
\end{reference}
Let $X$ be a scheme.
Let $\mathcal{L}$ be an invertible $\mathcal{O}_X$-module.
We say $\mathcal{L}$ is {\it ample} if
\begin{enumerate}
\item $X$ is quasi-compact, and
\item for every $x \in X$ there exists an $n \geq 1$
and $s \in \Gamma(X, \mathcal{L}^{\otimes n})$ such
that $x \in X_s$ and $X_s$ is affine.
\end{enumerate}
\end{definition}

\begin{definition}
\label{schemes-definition-reduced-induced-scheme}
Let $X$ be a scheme. Let $Z \subset X$ be a closed subset.
A {\it scheme structure on $Z$} is given by a closed subscheme $Z'$ of
$X$ whose underlying set is equal to $Z$. We often say
``let $(Z, \mathcal{O}_Z)$ be a scheme structure on $Z$'' to
indicate this. The {\it reduced induced scheme structure}
on $Z$ is the one constructed in Lemma \href{https://stacks.math.columbia.edu/tag/01J3}{lemma reduced closed subscheme (Tag 01J3)}.
The {\it reduction $X_{red}$ of $X$} is the reduced induced scheme
structure on $X$ itself.
\end{definition}

\begin{lemma}
\label{lemma-descend-opens}
In Situation \ref{situation-descent} we have the following:
\begin{enumerate}
\item Given any quasi-compact open $V \subset S = \lim_i S_i$
there exists an $i \in I$ and a quasi-compact open $V_i \subset S_i$
such that $f_i^{-1}(V_i) = V$.
\item Given $V_i \subset S_i$ and $V_{i'} \subset S_{i'}$
quasi-compact opens such that $f_i^{-1}(V_i) = f_{i'}^{-1}(V_{i'})$
there exists an index $i'' \geq i, i'$ such that
$f_{i''i}^{-1}(V_i) = f_{i''i'}^{-1}(V_{i'})$.
\item If $V_{1, i}, \ldots, V_{n, i} \subset S_i$ are quasi-compact
opens and $S = f_i^{-1}(V_{1, i}) \cup \ldots \cup f_i^{-1}(V_{n, i})$
then $S_{i'} = f_{i'i}^{-1}(V_{1, i}) \cup \ldots \cup f_{i'i}^{-1}(V_{n, i})$
for some $i' \geq i$.
\end{enumerate}
\end{lemma}

\begin{proof}
Choose $i_0 \in I$. Note that $I$ is nonempty as the limit is directed.
For convenience we write $S_0 = S_{i_0}$ and $i_0 = 0$.
Choose an affine open covering $S_0 = U_{1, 0} \cup \ldots \cup U_{m, 0}$.
Denote $U_{j, i} \subset S_i$ the inverse image of $U_{j, 0}$
under the transition morphism for $i \geq 0$.
Denote $U_j$ the inverse image of $U_{j, 0}$ in $S$.
Note that $U_j = \lim_i U_{j, i}$ is a limit of affine
schemes.

\medskip\noindent
We first prove the uniqueness statement: Let
$V_i \subset S_i$ and $V_{i'} \subset S_{i'}$
quasi-compact opens such that $f_i^{-1}(V_i) = f_{i'}^{-1}(V_{i'})$.
It suffices to show that $f_{i''i}^{-1}(V_i \cap U_{j, i''})$ and
$f_{i''i'}^{-1}(V_{i'} \cap U_{j, i''})$ become equal
for $i''$ large enough. Hence we reduce to the case
of a limit of affine schemes. In this case write
$S = \Spec(R)$ and $S_i = \Spec(R_i)$ for all $i \in I$.
We may write $V_i = S_i \setminus V(h_1, \ldots, h_m)$
and $V_{i'} = S_{i'} \setminus V(g_1, \ldots, g_n)$.
The assumption means that the ideals
$\sum g_jR$ and $\sum h_jR$ have the same radical
in $R$. This means that $g_j^N = \sum a_{jj'}h_{j'}$ and
$h_j^N = \sum b_{jj'} g_{j'}$ for some $N \gg 0$ and $a_{jj'}$
and $b_{jj'}$ in $R$.
Since $R = \colim_i R_i$ we can chose an index
$i'' \geq i$ such that the equations
$g_j^N = \sum a_{jj'}h_{j'}$ and
$h_j^N = \sum b_{jj'} g_{j'}$ hold in $R_{i''}$ for some
$a_{jj'}$ and $b_{jj'}$ in $R_{i''}$. This implies that
the ideals $\sum g_jR_{i''}$ and $\sum h_jR_{i''}$ have the same radical
in $R_{i''}$ as desired.

\medskip\noindent
We prove existence: If $S_0$ is affine, then $S_i = \Spec(R_i)$ for all
$i \geq 0$ and $S = \Spec(R)$ with $R = \colim R_i$. Then
$V = S \setminus V(g_1, \ldots, g_n)$ for some $g_1, \ldots, g_n \in R$.
Choose any $i$ large enough so that each of the $g_j$ comes from an
element $g_{j, i} \in R_i$ and take
$V_i = S_i \setminus V(g_{1, i}, \ldots, g_{n, i})$.
If $S_0$ is general, then the opens $V \cap U_j$
are quasi-compact because $S$ is quasi-separated. Hence by the
affine case we see that for each $j = 1, \ldots, m$
there exists an $i_j \in I$ and a quasi-compact open
$V_{i_j} \subset U_{j, i_j}$ whose inverse image in $U_j$
is $V \cap U_j$. Set $i = \max(i_1, \ldots, i_m)$
and let $V_i = \bigcup f_{ii_j}^{-1}(V_{i_j})$.

\medskip\noindent
The statement on coverings follows from the uniqueness statement
for the opens $V_{1, i} \cup \ldots \cup V_{n, i}$ and $S_i$ of $S_i$.
\end{proof}


\begin{lemma}
\label{lemma-limit-quasi-affine}
In Situation \ref{situation-descent} if $S$ is quasi-affine, then
for some $i_0 \in I$ the schemes $S_i$ for $i \geq i_0$ are quasi-affine.
\end{lemma}

\begin{proof}
Choose $i_0 \in I$. Note that $I$ is nonempty as the limit is directed.
For convenience we write $S_0 = S_{i_0}$ and $i_0 = 0$.
Let $s \in S$. We may choose an affine open
$U_0 \subset S_0$ containing $f_0(s)$. Since $S$ is quasi-affine
we may choose an element $a \in \Gamma(S, \mathcal{O}_S)$ such
that $s \in D(a) \subset f_0^{-1}(U_0)$, and such that
$D(a)$ is affine. By Lemma \ref{lemma-descend-section}
there exists an $i \geq 0$ such that $a$
comes from an element $a_i \in \Gamma(S_i, \mathcal{O}_{S_i})$.
For any index $j \geq i$ we denote $a_j$
the image of $a_i$ in the global sections of the
structure sheaf of $S_j$.
Consider the opens $D(a_j) \subset S_j$
and $U_j = f_{j0}^{-1}(U_0)$. Note that
$U_j$ is affine and $D(a_j)$ is a quasi-compact open of $S_j$,
see Properties, Lemma \href{https://stacks.math.columbia.edu/tag/01PV}{lemma affine cap s open (Tag 01PV)}
for example. Hence we may apply Lemma \ref{lemma-descend-opens} to the opens
$U_j$ and $U_j \cup D(a_j)$ to conclude that
$D(a_j) \subset U_j$ for some  $j \geq i$.
For such an index $j$ we see that $D(a_j) \subset S_j$ is an affine open
(because $D(a_j)$ is a standard affine open of the affine open $U_j$)
containing the image $f_j(s)$.

\medskip\noindent
We conclude that for every $s \in S$ there exist
an index $i \in I$, and a global section
$a \in \Gamma(S_i, \mathcal{O}_{S_i})$
such that $D(a) \subset S_i$ is an affine open
containing $f_i(s)$. Because $S$ is quasi-compact we
may choose a single index $i \in I$ and global sections
$a_1, \ldots, a_m \in \Gamma(S_i, \mathcal{O}_{S_i})$
such that each $D(a_j) \subset S_i$ is affine open
and such that $f_i : S \to S_i$ has image contained
in the union $W_i = \bigcup_{j = 1, \ldots, m} D(a_j)$.
For $i' \geq i$ set $W_{i'} = f_{i'i}^{-1}(W_i)$.
Since $f_i^{-1}(W_i)$ is all of $S$ we see
(by Lemma \ref{lemma-descend-opens} again)
that for a suitable $i' \geq i$ we
have $S_{i'} = W_{i'}$. Thus we may replace $i$ by
$i'$ and assume that $S_i = \bigcup_{j = 1, \ldots, m} D(a_j)$.
This implies that $\mathcal{O}_{S_i}$ is an ample invertible
sheaf on $S_i$ (see Properties, Definition \ref{properties-definition-ample})
and hence that $S_i$ is quasi-affine, see
Properties, Lemma \href{https://stacks.math.columbia.edu/tag/01QE}{lemma quasi affine O ample (Tag 01QE)}.
Hence we win.
\end{proof}


\begin{lemma}
\label{lemma-descend-section}
In Situation \ref{situation-descent}.
Suppose that $\mathcal{F}_0$ is a quasi-coherent sheaf on $S_0$.
Set $\mathcal{F}_i = f_{i0}^*\mathcal{F}_0$ for $i \geq 0$ and set
$\mathcal{F} = f_0^*\mathcal{F}_0$.
Then
$$
\Gamma(S, \mathcal{F}) = \colim_{i \geq 0} \Gamma(S_i, \mathcal{F}_i)
$$
\end{lemma}

\begin{proof}
Write $\mathcal{A}_j = f_{i0, *} \mathcal{O}_{S_i}$.
This is a quasi-coherent sheaf of $\mathcal{O}_{S_0}$-algebras
(see Morphisms, Lemma \href{https://stacks.math.columbia.edu/tag/01SA}{lemma affine equivalence algebras (Tag 01SA)})
and $S_i$ is the relative spectrum of $\mathcal{A}_i$ over $S_0$.
In the proof of Lemma \ref{lemma-directed-inverse-system-has-limit}
we constructed $S$ as the relative spectrum of
$\mathcal{A} = \colim_{i \geq 0} \mathcal{A}_i$
over $S_0$. Set
$$
\mathcal{M}_i = \mathcal{F}_0 \otimes_{\mathcal{O}_{S_0}} \mathcal{A}_i
$$
and
$$
\mathcal{M} = \mathcal{F}_0 \otimes_{\mathcal{O}_{S_0}} \mathcal{A}.
$$
Then we have $f_{i0, *} \mathcal{F}_i = \mathcal{M}_i$
and $f_{0, *}\mathcal{F} = \mathcal{M}$. Since $\mathcal{A}$
is the colimit of the sheaves $\mathcal{A}_i$ and since tensor
product commutes with directed colimits, we conclude that
$\mathcal{M} = \colim_{i \geq 0} \mathcal{M}_i$.
Since $S_0$ is quasi-compact and quasi-separated we see that
\begin{eqnarray*}
\Gamma(S, \mathcal{F})
& = &
\Gamma(S_0, \mathcal{M}) \\
& = &
\Gamma(S_0, \colim_{i \geq 0} \mathcal{M}_i) \\
& = &
\colim_{i \geq 0} \Gamma(S_0, \mathcal{M}_i) \\
& = &
\colim_{i \geq 0} \Gamma(S_i, \mathcal{F}_i)
\end{eqnarray*}
see Sheaves, Lemma \href{https://stacks.math.columbia.edu/tag/009F}{lemma directed colimits sections (Tag 009F)} and
Topology, Lemma \href{https://stacks.math.columbia.edu/tag/0069}{lemma topology quasi separated scheme (Tag 0069)}
for the middle equality.
\end{proof}


\begin{lemma}
\label{lemma-directed-inverse-system-has-limit}
Let $I$ be a directed set. Let $(S_i, f_{ii'})$ be an
inverse system of schemes over $I$. If all the morphisms
$f_{ii'} : S_i \to S_{i'}$ are affine, then the limit $S = \lim_i S_i$ exists
in the category of schemes. Moreover,
\begin{enumerate}
\item each of the morphisms $f_i : S \to S_i$ is affine,
\item for an element $0 \in I$ and any open subscheme $U_0 \subset S_0$
we have
$$
f_0^{-1}(U_0) = \lim_{i \geq 0} f_{i0}^{-1}(U_0)
$$
in the category of schemes.
\end{enumerate}
\end{lemma}

\begin{proof}
Choose an element $0 \in I$. Note that $I$ is nonempty as the limit is
directed. For every $i \geq 0$ consider the quasi-coherent sheaf of
$\mathcal{O}_{S_0}$-algebras $\mathcal{A}_i = f_{i0, *}\mathcal{O}_{S_i}$.
Recall that $S_i = \underline{\Spec}_{S_0}(\mathcal{A}_i)$,
see Morphisms, Lemma \href{https://stacks.math.columbia.edu/tag/01S8}{lemma characterize affine (Tag 01S8)}.
Set $\mathcal{A} = \colim_{i \geq 0} \mathcal{A}_i$.
This is a quasi-coherent sheaf of $\mathcal{O}_{S_0}$-algebras,
see Schemes, Section \href{https://stacks.math.columbia.edu/tag/01LA}{section quasi coherent (Tag 01LA)}.
Set $S = \underline{\Spec}_{S_0}(\mathcal{A})$.
By Morphisms, Lemma \href{https://stacks.math.columbia.edu/tag/01SA}{lemma affine equivalence algebras (Tag 01SA)}
we get for $i \geq 0$ morphisms $f_i : S \to S_i$ compatible with
the transition morphisms. Note that the morphisms $f_i$ are
affine by Morphisms, Lemma \href{https://stacks.math.columbia.edu/tag/01SG}{lemma affine permanence (Tag 01SG)} for example.
By Lemma \href{https://stacks.math.columbia.edu/tag/01YW}{lemma directed inverse system affine schemes has limit (Tag 01YW)} above
we see that for any affine open $U_0 \subset S_0$ the
inverse image $U = f_0^{-1}(U_0) \subset S$ is the limit of the
system of opens $U_i = f_{i0}^{-1}(U_0)$, $i \geq 0$ in the
category of schemes.

\medskip\noindent
Let $T$ be a scheme. Let $g_i : T \to S_i$ be a compatible system
of morphisms. To show that $S = \lim_i S_i$ we have
to prove there is a unique morphism $g : T \to S$ with
$g_i = f_i \circ g$ for all $i \in I$.
For every $t \in T$ there exists an affine open
$U_0 \subset S_0$ containing $g_0(t)$. Let $V \subset g_0^{-1}(U_0)$
be an affine open neighbourhood containing $t$.
By the remarks above we obtain a unique morphism
$g_V : V \to U = f_0^{-1}(U_0)$ such that $f_i \circ g_V = g_i|_{U_i}$
for all $i$. The open sets $V \subset T$ so constructed form
a basis for the topology of $T$. The morphisms $g_V$ glue to a morphism
$g : T \to S$ because of the uniqueness property. This gives the
desired morphism $g : T \to S$.

\medskip\noindent
The final statement is clear from the construction of the limit above.
\end{proof}


\begin{lemma}
\label{lemma-limit-affine}
In Situation \ref{situation-descent} if $S$ is affine,
then for some $i_0 \in I$ the schemes $S_i$ for $i \geq i_0$
are affine.
\end{lemma}

\begin{proof}
By Lemma \ref{lemma-limit-quasi-affine} we may assume that $S_0$ is
quasi-affine for some $0 \in I$. Set $R_0 = \Gamma(S_0, \mathcal{O}_{S_0})$.
Then $S_0$ is a quasi-compact open of $T_0 = \Spec(R_0)$. Denote
$j_0 : S_0 \to T_0$ the corresponding quasi-compact open immersion.
For $i \geq 0$ set $\mathcal{A}_i = f_{i0, *}\mathcal{O}_{S_i}$.
Since $f_{i0}$ is affine we see that
$S_i = \underline{\Spec}_{S_0}(\mathcal{A}_i)$.
Set $T_i = \underline{\Spec}_{T_0}(j_{0, *}\mathcal{A}_i)$.
Then $T_i \to T_0$ is affine, hence $T_i$ is affine. Thus
$T_i$ is the spectrum of
$$
R_i = \Gamma(T_0, j_{0, *}\mathcal{A}_i) = \Gamma(S_0, \mathcal{A}_i) =
\Gamma(S_i, \mathcal{O}_{S_i}).
$$
Write $S = \Spec(R)$. We have $R = \colim_i R_i$
by Lemma \ref{lemma-descend-section}.
Hence also $S = \lim_i T_i$. As formation of the relative spectrum commutes
with base change, the inverse image
of the open $S_0 \subset T_0$ in $T_i$ is $S_i$.
Let $Z_0 = T_0 \setminus S_0$ and let $Z_i \subset T_i$
be the inverse image of $Z_0$. As $S_i = T_i \setminus Z_i$, it suffices
to show that $Z_i$ is empty for some $i$. Assume $Z_i$ is nonempty for all
$i$ to get a contradiction. By Lemma \ref{lemma-limit-closed-nonempty}
there exists a point $s$ of $S = \lim T_i$ which maps to a point of $Z_i$
for every $i$. But $S = \lim_i S_i$, and hence we arrive at a contradiction
by Lemma \ref{lemma-topology-limit}.
\end{proof}


\begin{lemma}
\label{lemma-limit-closed-nonempty}
In Situation \ref{situation-descent}.
Suppose for each $i$ we are given a nonempty closed subset
$Z_i \subset S_i$ with $f_{i'i}(Z_{i'}) \subset Z_i$ for all
$i' \geq i$.
Then there exists a point $s \in S$ with $f_i(s) \in Z_i$ for
all $i$.
\end{lemma}

\begin{proof}
Let $Z_i \subset S_i$ also denote the reduced closed subscheme
associated to $Z_i$, see Schemes,
Definition \ref{schemes-definition-reduced-induced-scheme}.
A closed immersion is affine, and a composition of affine
morphisms is affine (see
Morphisms, Lemmas \href{https://stacks.math.columbia.edu/tag/01SE}{lemma closed immersion affine (Tag 01SE)}
and \href{https://stacks.math.columbia.edu/tag/01SC}{lemma composition affine (Tag 01SC)}), and hence $Z_{i'} \to S_i$ is
affine when $i' \geq i$. We conclude that the morphism
$f_{i'i} : Z_{i'} \to Z_i$ is affine by
Morphisms, Lemma \href{https://stacks.math.columbia.edu/tag/01SG}{lemma affine permanence (Tag 01SG)}.
Each of the schemes $Z_i$ is quasi-compact as a closed
subscheme of a quasi-compact scheme. Hence we may apply
Lemma \ref{lemma-limit-nonempty} to see that
$Z = \lim_i Z_i$ is nonempty. Since there is a
canonical morphism $Z \to S$ we win.
\end{proof}


\begin{lemma}
\label{lemma-limit-nonempty}
Let $S = \lim S_i$ be the limit of a directed inverse system
of schemes with affine transition morphisms
(Lemma \ref{lemma-directed-inverse-system-has-limit}).
If all the schemes $S_i$ are nonempty and quasi-compact,
then the limit $S = \lim_i S_i$ is nonempty.
\end{lemma}

\begin{proof}
Choose $0 \in I$. Note that $I$ is nonempty as the limit is directed.
Choose an affine open covering $S_0 = \bigcup_{j = 1, \ldots, m} U_j$.
Since $I$ is directed there exists a $j \in \{1, \ldots, m\}$
such that $f_{i0}^{-1}(U_j) \not = \emptyset$ for all
$i \geq 0$. Hence $\lim_{i \geq 0} f_{i0}^{-1}(U_j)$ is not
empty since a directed colimit of nonzero rings is nonzero
(because $1 \not = 0$). As $\lim_{i \geq 0} f_{i0}^{-1}(U_j)$
is an open subscheme of the limit we win.
\end{proof}


\begin{lemma}
\label{lemma-topology-limit}
In Situation \ref{situation-descent}.
\begin{enumerate}
\item We have $S_{set} = \lim_i S_{i, set}$ where $S_{set}$
indicates the underlying set of the scheme $S$.
\item We have $S_{top} = \lim_i S_{i, top}$ where $S_{top}$
indicates the underlying topological space of the scheme $S$.
\item If $s, s' \in S$ and $s'$ is not a specialization of $s$
then for some $i \in I$ the image $s'_i \in S_i$ of $s'$ is not
a specialization of the image $s_i \in S_i$ of $s$.
\item Add more easy facts on topology of $S$ here.
(Requirement: whatever is added should be easy in the affine case.)
\end{enumerate}
\end{lemma}

\begin{proof}
Part (1) is a special case of Lemma \href{https://stacks.math.columbia.edu/tag/0CUE}{lemma inverse limit sets (Tag 0CUE)}.

\medskip\noindent
Part (2) is a special case of Lemma \href{https://stacks.math.columbia.edu/tag/0CUF}{lemma inverse limit top (Tag 0CUF)}.

\medskip\noindent
Part (3) is a special case of Lemma \href{https://stacks.math.columbia.edu/tag/0CUG}{lemma inverse limit irreducibles (Tag 0CUG)}.
\end{proof}


\begin{lemma}
\label{lemma-scheme-over-limit}
Let $I$ be a directed set.
Let $(S_i, f_{ii'})$ be an inverse system of schemes over $I$.
Assume all the morphisms $f_{ii'} : S_i \to S_{i'}$ are affine,
Let $S = \lim_i S_i$. Let $0 \in I$.
Suppose that $T$ is a scheme over $S_0$.
Then
$$
T \times_{S_0} S = \lim_{i \geq 0} T \times_{S_0} S_i
$$
\end{lemma}

\begin{proof}
The right hand side is a scheme by
Lemma \ref{lemma-directed-inverse-system-has-limit}.
The equality is formal, see
Categories, Lemma \href{https://stacks.math.columbia.edu/tag/002M}{lemma colimits commute (Tag 002M)}.
\end{proof}


\begin{lemma}
\label{algebra-lemma-colimit-category-fp-algebras}
Suppose that $R = \colim_{\lambda \in \Lambda} R_\lambda$ is a directed colimit
of rings. Then the category of finitely presented $R$-algebras is
the colimit of the categories of finitely presented $R_\lambda$-algebras.
More precisely
\begin{enumerate}
\item Given a finitely presented $R$-algebra $A$ there exists a
$\lambda \in \Lambda$ and a finitely presented $R_\lambda$-algebra
$A_\lambda$ such that $A \cong A_\lambda \otimes_{R_\lambda} R$.
\item Given a $\lambda \in \Lambda$, finitely presented
$R_\lambda$-algebras $A_\lambda, B_\lambda$, and an $R$-algebra map
$\varphi : A_\lambda \otimes_{R_\lambda} R \to B_\lambda \otimes_{R_\lambda} R$,
then there exists a $\mu \geq \lambda$ and an $R_\mu$-algebra map
$\varphi_\mu : A_\lambda \otimes_{R_\lambda} R_\mu \to
B_\lambda \otimes_{R_\lambda} R_\mu$
such that $\varphi = \varphi_\mu \otimes 1_R$.
\item Given a $\lambda \in \Lambda$, finitely presented
$R_\lambda$-algebras $A_\lambda, B_\lambda$, and $R_\lambda$-algebra maps
$\varphi_\lambda, \psi_\lambda : A_\lambda \to B_\lambda$
such that $\varphi \otimes 1_R = \psi \otimes 1_R$, then
$\varphi \otimes 1_{R_\mu} = \psi \otimes 1_{R_\mu}$ for some
$\mu \geq \lambda$.
\end{enumerate}
\end{lemma}

\begin{proof}
To prove (1) choose a presentation
$A = R[x_1, \ldots, x_n]/(f_1, \ldots, f_m)$.
We can choose a $\lambda \in \Lambda$ and elements
$f_{\lambda, j} \in R_\lambda[x_1, \ldots, x_n]$ mapping to
$f_j \in R[x_1, \ldots, x_n]$.
Then we simply let
$A_\lambda =
R_\lambda[x_1, \ldots, x_n]/(f_{\lambda, 1}, \ldots, f_{\lambda, m})$.

\medskip\noindent
Parts (2) and (3) follow from
Lemma \ref{algebra-lemma-algebra-map-property-in-colimit}.
\end{proof}


\begin{lemma}
\label{algebra-lemma-algebra-map-property-in-colimit}
Let $A$ be a ring and let $B, C$ be $A$-algebras.
Suppose that $R = \colim_{i \in I} R_i$ is a directed colimit
of $A$-algebras.
\begin{enumerate}
\item If $B$ is a finite type $A$-algebra, and $u, u' : B \to C$ are
$A$-algebra maps such that
$u \otimes 1 = u' \otimes 1 : B \otimes_A R \to C \otimes_A R$
then for some $i$ we have
$u \otimes 1 = u' \otimes 1 : B \otimes_A R_i \to C \otimes_A R_i$.
\item If $C$ is a finite type $A$-algebra and $u : B \to C$ is an
$A$-algebra map such that
$u \otimes 1 : B \otimes_A R \to C \otimes_A R$ is surjective, then
for some $i$ the map $u \otimes 1 : B \otimes_A R_i \to C \otimes_A R_i$
is surjective.
\item If $C$ is of finite presentation over $A$ and
$v : C \otimes_A R \to B \otimes_A R$ is an $R$-algebra map, then there
exists an $i$ and an $R_i$-algebra map
$v_i : C \otimes_A R_i \to B \otimes_A R_i$ such that
$v = v_i \otimes 1$.
\item If $B$ is a finite type $A$-algebra, $C$ is a finitely presented
$A$-algebra, and
$u \otimes 1 : B \otimes_A R \to C \otimes_A R$ is an isomorphism, then
for some $i$ the map $u \otimes 1 : B \otimes_A R_i \to C \otimes_A R_i$
is an isomorphism.
\end{enumerate}
\end{lemma}

\begin{proof}
To prove (1) assume $u$ is as in (1) and
let $x_1, \ldots, x_m \in B$ be generators. Since
$B \otimes_A R = \colim_i B \otimes_A R_i$
we may pick an $i \in I$ such that $u(x_j) \otimes 1 = u'(x_j) \otimes 1$
in $B \otimes_A R_i$, $j = 1, \ldots, m$.
For such an $i$ we have
$u \otimes 1 = u' \otimes 1 : B \otimes_A R_i \to C \otimes_A R_i$.

\medskip\noindent
To prove (2) assume $u \otimes 1$ surjective and
let $y_1, \ldots, y_m \in C$ be generators. Since
$B \otimes_A R = \colim_i B \otimes_A R_i$
we may pick an $i \in I$ and $z_j \in B \otimes_A R_i$, $j = 1, \ldots, m$
whose images in $C \otimes_A R$ equal $y_j \otimes 1$.
For such an $i$ the map $u \otimes 1 : B \otimes_A R_i \to C \otimes_A R_i$
is surjective.

\medskip\noindent
To prove (3) let $c_1, \ldots, c_m \in C$ be generators. Let
$K = \Ker(A[x_1, \ldots, x_m] \to N)$ where the map is given by
the rule $x_j \mapsto \sum c_j$. Let $f_1, \ldots, f_t$
be generators for $K$ as an ideal in $A[x_1, \ldots, x_m]$.
We think of $f_j = f_j(x_1, \ldots, x_m)$ as a polynomial.
Since $B \otimes_A R = \colim_i B \otimes_A R_i$
we may pick an $i \in I$ and $z_j \in B \otimes_A R_i$, $j = 1, \ldots, m$
whose images in $B \otimes_A R$ equal $v(c_j \otimes 1)$.
We want to use the $z_j$ to define a map
$v_i : C \otimes_A R_i \to B \otimes_A R_i$.
Since $K \otimes_A R_i \to R_i[x_1, \ldots, x_m] \to C \otimes_A R_i \to 0$
is a presentation, it suffices to check that
$\xi_s = f_j(z_1, \ldots, z_m)$ is
zero in $B \otimes_A R_i$ for each $s = 1, \ldots, t$. This may not
be the case, but since the image of $\xi_s$ in $B \otimes_A R$ is zero
we see that it will be the case after increasing $i$ a bit.

\medskip\noindent
To prove (4) assume $u \otimes 1$ is an isomorphism, that
$B$ is a finite type $A$-algebra, and that $C$ is a finitely presented
$A$-algebra. Let $v : B \otimes_A R \to C \otimes_A R$ be an inverse to
$u \otimes 1$. Let $v_i : C \otimes_A R_i \to B \otimes_A R_i$ be as
in part (3). Apply part (1) to see that, after increasing $i$ we have
$v_i \circ (u \otimes 1) = \text{id}_{B \otimes_R R_i}$ and
$(u \otimes 1) \circ v_i = \text{id}_{C \otimes_R R_i}$.
\end{proof}


\begin{lemma}
\label{lemma-descend-finite-presentation}
Let $I$ be a directed set.
Let $(S_i, f_{ii'})$ be an inverse system of schemes over $I$.
Assume
\begin{enumerate}
\item the morphisms $f_{ii'} : S_i \to S_{i'}$ are affine,
\item the schemes $S_i$ are quasi-compact and quasi-separated.
\end{enumerate}
Let $S = \lim_i S_i$. Then we have the following:
\begin{enumerate}
\item For any morphism of finite presentation $X \to S$
there exists an index $i \in I$ and a morphism of finite
presentation $X_i \to S_i$ such that $X \cong X_{i, S}$ as
schemes over $S$.
\item Given an index $i \in I$, schemes
$X_i$, $Y_i$ of finite presentation over $S_i$, and a morphism
$\varphi : X_{i, S} \to Y_{i, S}$ over $S$, there exists an index
$i' \geq i$ and a morphism
$\varphi_{i'} : X_{i, S_{i'}} \to Y_{i, S_{i'}}$
whose base change to $S$ is $\varphi$.
\item Given an index $i \in I$, schemes $X_i$, $Y_i$ of finite presentation
over $S_i$ and a pair of morphisms $\varphi_i, \psi_i : X_i \to Y_i$
whose base changes $\varphi_{i, S} = \psi_{i, S}$ are equal,
there exists an index $i' \geq i$ such that
$\varphi_{i, S_{i'}} = \psi_{i, S_{i'}}$.
\end{enumerate}
In other words, the category of schemes of finite presentation over
$S$ is the colimit over $I$ of the categories of schemes of finite
presentation over $S_i$.
\end{lemma}

\begin{proof}
In case each of the schemes $S_i$ is affine, and we consider
only affine schemes of finite presentation over $S_i$, resp.\ $S$
this lemma is equivalent to
Algebra, Lemma \ref{algebra-lemma-colimit-category-fp-algebras}.
We claim that the affine case implies the lemma in general.

\medskip\noindent
Let us prove (3). Suppose given an index $i \in I$, schemes
$X_i$, $Y_i$ of finite presentation over $S_i$ and a pair of morphisms
$\varphi_i, \psi_i : X_i \to Y_i$. Assume that the base changes are
equal: $\varphi_{i, S} = \psi_{i, S}$. We will use the notation
$X_{i'} = X_{i, S_{i'}}$ and $Y_{i'} = Y_{i, S_{i'}}$ for
$i' \geq i$. We also set $X = X_{i, S}$ and $Y = Y_{i, S}$.
Note that according to Lemma \ref{lemma-scheme-over-limit} we have
$X = \lim_{i' \geq i} X_{i'}$ and similarly for $Y$.
Additionally we denote $\varphi_{i'}$ and $\psi_{i'}$
(resp.\ $\varphi$ and $\psi$)
the base change of $\varphi_i$ and $\psi_i$ to $S_{i'}$
(resp.\ $S$). So our assumption means that $\varphi = \psi$.
Since $Y_i$ and $X_i$ are of finite presentation
over $S_i$, and since $S_i$ is quasi-compact and quasi-separated, also
$X_i$ and $Y_i$ are quasi-compact and quasi-separated
(see Morphisms,
Lemma \href{https://stacks.math.columbia.edu/tag/01TY}{lemma finite presentation quasi compact quasi separated (Tag 01TY)}).
Hence we may choose a finite affine open covering
$Y_i = \bigcup V_{j, i}$ such that each $V_{j, i}$ maps into
an affine open of $S_i$. As above, denote $V_{j, i'}$ the inverse
image of $V_{j, i}$ in $Y_{i'}$ and $V_j$ the inverse image in $Y$.
The immersions $V_{j, i'} \to Y_{i'}$ are quasi-compact, and the inverse images
$U_{j, i'} = \varphi_i^{-1}(V_{j, i'})$ and
$U_{j, i'}' = \psi_i^{-1}(V_{j, i'})$
are quasi-compact opens of $X_{i'}$. By assumption the inverse images of
$V_j$ under $\varphi$ and $\psi$ in $X$ are equal.
Hence by Lemma \ref{lemma-descend-opens}
there exists an index $i' \geq i$ such that
of $U_{j, i'} = U_{j, i'}'$ in $X_{i'}$.
Choose an finite affine open covering
$U_{j, i'} = U_{j, i'}' = \bigcup W_{j, k, i'}$
which induce coverings $U_{j, i''} = U_{j, i''}' = \bigcup W_{j, k, i''}$
for all $i'' \geq i'$.
By the affine case there exists
an index $i''$ such that
$\varphi_{i''}|_{W_{j, k, i''}} = \psi_{i''}|_{W_{j, k, i''}}$
for all $j, k$. Then $i''$ is an index such that
$\varphi_{i''} = \psi_{i''}$ and (3) is proved.

\medskip\noindent
Let us prove (2). Suppose given an index $i \in I$, schemes
$X_i$, $Y_i$ of finite presentation over $S_i$ and a morphism
$\varphi : X_{i, S} \to Y_{i, S}$. We will use the notation
$X_{i'} = X_{i, S_{i'}}$ and $Y_{i'} = Y_{i, S_{i'}}$ for
$i' \geq i$. We also set $X = X_{i, S}$ and $Y = Y_{i, S}$.
Note that according to Lemma \ref{lemma-scheme-over-limit} we have
$X = \lim_{i' \geq i} X_{i'}$ and similarly for $Y$.
Since $Y_i$ and $X_i$ are of finite presentation
over $S_i$, and since $S_i$ is quasi-compact and quasi-separated, also
$X_i$ and $Y_i$ are quasi-compact and quasi-separated
(see Morphisms,
Lemma \href{https://stacks.math.columbia.edu/tag/01TY}{lemma finite presentation quasi compact quasi separated (Tag 01TY)}).
Hence we may choose a finite affine open covering
$Y_i = \bigcup V_{j, i}$ such that each $V_{j, i}$ maps into
an affine open of $S_i$. As above, denote $V_{j, i'}$ the inverse
image of $V_{j, i}$ in $Y_{i'}$ and $V_j$ the inverse image in $Y$.
The immersions $V_j \to Y$ are quasi-compact, and the inverse images
$U_j = \varphi^{-1}(V_j)$ are quasi-compact opens of $X$.
Hence by Lemma \ref{lemma-descend-opens} there exists an index
$i' \geq i$ and quasi-compact opens $U_{j, i'}$ of $X_{i'}$
whose inverse image in $X$ is $U_j$. Choose an finite affine open covering
$U_{j, i'} = \bigcup W_{j, k, i'}$ which induce affine open
coverings $U_{j, i''} = \bigcup W_{j, k, i''}$
for all $i'' \geq i'$ and an affine open covering
$U_j = \bigcup W_{j, k}$. By the affine case there exists
an index $i''$ and morphisms
$\varphi_{j, k, i''} : W_{j, k, i''} \to V_{j, i''}$
such that
$\varphi|_{W_{j, k}} = \varphi_{j, k, i'', S}$ for all $j, k$.
By part (3) proved above, there is a further index $i''' \geq i''$
such that
$$
\varphi_{j_1, k_1, i'', S_{i'''}}|_{W_{j_1, k_1, i'''} \cap W_{j_2, k_2, i'''}}
=
\varphi_{j_2, k_2, i'', S_{i'''}}|_{W_{j_1, k_1, i'''} \cap W_{j_2, k_2, i'''}}
$$
for all $j_1, j_2, k_1, k_2$. Then $i'''$ is an index such that
there exists a morphism $\varphi_{i'''} : X_{i'''} \to Y_{i'''}$
whose base change to $S$ gives $\varphi$. Hence (2) holds.

\medskip\noindent
Let us prove (1). Suppose given a scheme $X$ of finite presentation
over $S$. Since $X$ is of finite presentation
over $S$, and since $S$ is quasi-compact and quasi-separated, also
$X$ is quasi-compact and quasi-separated
(see Morphisms,
Lemma \href{https://stacks.math.columbia.edu/tag/01TY}{lemma finite presentation quasi compact quasi separated (Tag 01TY)}).
Choose a finite affine open covering $X = \bigcup U_j$
such that each $U_j$ maps into an affine open $V_j \subset S$.
Denote $U_{j_1j_2} = U_{j_1} \cap U_{j_2}$ and
$U_{j_1j_2j_3} = U_{j_1} \cap U_{j_2} \cap U_{j_3}$.
By Lemmas \ref{lemma-descend-opens} and \ref{lemma-limit-affine}
we can find an index $i_1$ and affine opens $V_{j, i_1} \subset S_{i_1}$
such that each $V_j$ is the inverse of this in $S$.
Let $V_{j, i}$ be the inverse image of $V_{j, i_1}$ in $S_i$ for
$i \geq i_1$. By the affine case we may find an index $i_2 \geq i_1$ and
affine schemes $U_{j, i_2} \to V_{j, i_2}$ such
that $U_j = S \times_{S_{i_2}} U_{j, i_2}$ is the base change.
Denote $U_{j, i} = S_i \times_{S_{i_2}} U_{j, i_2}$ for $i \geq i_2$.
By Lemma \ref{lemma-descend-opens} there exists an index
$i_3 \geq i_2$ and open subschemes
$W_{j_1, j_2, i_3} \subset U_{j_1, i_3}$
whose base change to $S$ is equal to $U_{j_1j_2}$.
Denote $W_{j_1, j_2, i} = S_i \times_{S_{i_3}} W_{j_1, j_2, i_3}$
for $i \geq i_3$. By part (2) shown above there exists an index
$i_4 \geq i_3$ and morphisms
$\varphi_{j_1, j_2, i_4} : W_{j_1, j_2, i_4} \to W_{j_2, j_1, i_4}$
whose base change to $S$ gives the identity morphism
$U_{j_1j_2} = U_{j_2j_1}$ for all $j_1, j_2$.
For all $i \geq i_4$ denote
$\varphi_{j_1, j_2, i} = \text{id}_S \times \varphi_{j_1, j_2, i_4}$
the base change. We claim that for some $i_5 \geq i_4$ the system
$((U_{j, i_5})_j, (W_{j_1, j_2, i_5})_{j_1, j_2},
(\varphi_{j_1, j_2, i_5})_{j_1, j_2})$ forms a glueing datum
as in Schemes, Section \href{https://stacks.math.columbia.edu/tag/01JA}{section glueing schemes (Tag 01JA)}.
In order to see this we have to verify that for $i$ large enough
we have
$$
\varphi_{j_1, j_2, i}^{-1}(W_{j_2, j_1, i} \cap W_{j_2, j_3, i})
=
W_{j_1, j_2, i} \cap W_{j_1, j_3, i}
$$
and that for large enough $i$ the cocycle condition holds.
The first condition follows from Lemma \ref{lemma-descend-opens}
and the fact that $U_{j_2j_1j_3} = U_{j_1j_2j_3}$.
The second from part (3) of the lemma proved above and the fact
that the cocycle condition holds for the maps
$\text{id} : U_{j_1j_2} \to U_{j_2j_1}$.
Ok, so now we can use Schemes, Lemma \href{https://stacks.math.columbia.edu/tag/01JC}{lemma glue schemes (Tag 01JC)}
to glue the system
$((U_{j, i_5})_j, (W_{j_1, j_2, i_5})_{j_1, j_2},
(\varphi_{j_1, j_2, i_5})_{j_1, j_2})$ to get a scheme
$X_{i_5} \to S_{i_5}$. By construction the base change of
$X_{i_5}$ to $S$ is formed by glueing the open affines
$U_j$ along the opens $U_{j_1} \leftarrow U_{j_1j_2} \rightarrow U_{j_2}$.
Hence $S \times_{S_{i_5}} X_{i_5} \cong X$ as desired.
\end{proof}
\end{document}
